% Section 5.6, from lectures/L05/README.md and appendix A.4.

\section{Summary}
\label{c5:sec:review}

\needspace{6\baselineskip}
\subsection*{What you should be able to do now}
\begin{itemize}
  \item Say what separates a \code{signal} from a \code{variable}, predict the different result
    each gives from the same code, and choose between them deliberately.
  \item Explain what an FPGA physically consists of, why two differently written descriptions of
    the same circuit give the same hardware, and why a Karnaugh map buys less here than on discrete
    logic.
  \item Explain what limits a design's clock frequency, and why more registers can make a design
    faster rather than slower.
  \item Recognise an unintended latch from its warning, say what caused it and fix it, and say why
    ``it compiled without errors'' promises less for an FPGA design than for a C program.
\end{itemize}

\needspace{6\baselineskip}
\subsection*{Questions to test yourself with}
\begin{enumerate}
  \item A \code{signal} is assigned twice, with different values, in one round through a
    \code{process}. What ends up on it, and when does that update actually happen?
  \item The same question for a \code{variable}.
  \item Why does accumulating over a loop with a \code{signal} give a different, usually incorrect,
    result?
  \item Nothing on an FPGA becomes a gate. What does a Boolean function become instead, and why
    does a function of four variables cost the same whether or not you minimised it first?
  \item What is the critical path, and why can more flip-flops make a design run faster?
  \item A combinational \code{case} assigns its output in three branches out of four. What does the
    synthesis tool build for the fourth, and why is that the only thing it \emph{can} build?
  \item Why is an unintended latch a warning and not an error, and why does that make it more
    dangerous rather than less?
\end{enumerate}

\needspace{6\baselineskip}
\subsection*{Looking ahead}
Nothing in \chapref{c6:ch} to \ref{c8:ch} introduces a new idiom. Counters, shift registers and
timers all hold state across clock edges using \chapref{c3:ch}'s clocked process and the scheduling
rule above, and \chapref{c8:ch}'s state machines merely add an enumerated type and a \code{case} on
it. What changes is what the state is used for.

Two things to watch for. \Chapref{c6:ch}'s \code{serial_rx8} is the first design whose correctness
you cannot check by eye on an LED, which is where the testbenches handed out stop being a
formality. And \secref{c4:sec:generics}'s generic is put to work straight away: \chapref{c7:ch}'s
\code{timer} uses one for its tick count, and \code{walking_led} instantiates that timer alongside
\code{reset_sync} and \code{button_sync}, writing almost no logic of its own.
